\FloatBarrier
\newpage
\section{积分的应用}
\label{sec:calculus-integral-applications}

\begingroup
\setlength{\parskip}{3pt}
\setlength{\abovedisplayskip}{10pt}
\setlength{\belowdisplayskip}{10pt}

积分应用的共同任务，是确定每小段贡献什么、沿哪个变量累积，再把局部贡献相加。不同问题的表达式来自数量关系，累积过程则由同一种积分运算完成。

\subsection{平面图形的面积}

设连续曲线\(y=f(x)\)、\(y=g(x)\)在\([a,b]\)上满足\(f\ge g\)。第\(i\)条窄带的宽度为\(\Delta x_i\)，高度可以近似为\(f(\xi_i)-g(\xi_i)\)，故区域面积由矩形和的极限给出：
\begin{equation}\label{eq:calculus-area-between-graphs}
 S=\int_a^b(f(x)-g(x))\,dx.
\end{equation}
若上下关系改变，则先求交点，再分段；也可以统一写为\(\int_a^b|f-g|\)。若区域更适合按水平方向切片，就改用\(y\)为积分变量，以右边界横坐标减去左边界横坐标作为窄带长度。

例如，\(y=x\)与\(y=x^2\)相交于\(x=0,1\)，在中间直线位于抛物线上方，所以
\begin{equation*}
 S=\int_0^1(x-x^2)\,dx=\frac16.
\end{equation*}

\begin{figure}[!ht]
\centering
\includegraphics[width=.44\linewidth]{calculus-area-slices}
\caption{两条曲线之间的竖直窄带：局部面积由高度差乘横向宽度近似。}
\label{fig:calculus-area-slices}
\end{figure}

切片方向影响计算的简便程度。若竖直切片截出几段，改用水平切片后可能只剩一段。因此应根据区域的形状选择积分变量。

\begin{example}
求\(y=x^2\)与\(y=2-x^2\)围成的面积。再求椭圆\(x^2/a^2+y^2/b^2=1\)内部的面积，利用\(x=au\)的换元与第四章伸缩变换的面积倍率相互核对。
另求\(y=|x|\)与\(y=x^2\)围成的总面积，说明积分区间怎样由交点确定。
\end{example}

\begin{exercise}
\mbox{}\par\nobreak

（1）求\(y=x^3\)与\(y=x\)围成的全部有限区域的总面积，说明为什么不能只计算\(\int_{-1}^1(x-x^3)\,dx\)。

（2）求\(x=y^2\)与\(x=2-y\)围成的面积，比较水平与竖直切片的计算。

（3）对\(0<a<1\)，求\(y=x^2\)、横轴及直线\(x=a\)、\(x=1\)围成区域的面积，并求\(a\)，使它恰为\([0,1]\)上对应总面积的一半。

（4）求\(y=\ln x\)、横轴及直线\(x=e\)围成的面积。再求\(y=e^x\)、纵轴及直线\(y=e\)围成的面积，比较两区域关于\(y=x\)的对称关系。
\end{exercise}

\subsection{总变化量}

若位置函数\(s\)连续可导，速度为\(v=s'\)，则基本定理给出
\begin{equation*}
 s(b)-s(a)=\int_a^b v(t)\,dt.
\end{equation*}
这是位移。在不改变运动方向的时间段，路程等于位移的绝对值；各段相加，得到
\begin{equation*}
 \text{总路程}=\int_a^b|v(t)|\,dt.
\end{equation*}
速度变号的位置决定怎样拆分，而积分的上下限由所研究的时间范围决定。

位移累积各段有正负的变化，路程累积各段经过的长度。

同样，若\(q(t)\)表示每单位时间流入的量，则\(q(t)\Delta t\)近似一小段时间的流入量，\(\int_a^bq(t)\,dt\)给出总流入量。若同时有流出，则净变化应使用流入率减去流出率。

对细杆，设\(\rho(x)\)表示每单位长度的质量，称为线密度。短段质量近似为\(\rho(x)\Delta x\)，故整根杆的质量为\(\int_a^b\rho(x)\,dx\)。被积量与微分长度的单位相乘，应与所求总量一致。

若物体沿直线运动，沿运动坐标方向的力为\(F(x)\)，恒力在位移\(\Delta x\)上做功\(F\Delta x\)。变力可以在短段内用一个取样值近似，所以从\(a\)运动到\(b\)的功为
\begin{equation*}
 W=\int_a^bF(x)\,dx.
\end{equation*}
例如，弹簧的拉力为\(F(x)=kx\)，其中\(x\)是伸长量，功由各段的力乘位移累积。

\begin{example}
质点速度为\(v(t)=t^2-2t\)，\(0\le t\le3\)，初始位置为\(s(0)=1\)。求位置函数、总位移和总路程。再设一个容器的流入率为\(2+t\)、流出率为\(t^2\)，求\([0,2]\)内储量的净变化，并判断储量何时最大。
\end{example}

若位置为\(x=x(t)\)，由换元公式，功也可写成\(\int F(x(t))x'(t)\,dt\)。因子\(x'(t)\)将时间增量转为有方向的位移增量。

\begin{example}
一根细杆位于\([0,L]\)，从左端到位置\(x\)的累计质量为\(M(x)=kx^2\)，其中\(k>0\)。求线密度及区间\([L/2,L]\)内的质量，用密度积分核对。累计质量与线密度记录了什么？
\end{example}

\subsection{加权平均}

将不同部分赋予不同的非负权重，可以研究整体的平均位置或平均水平。设\(w\)连续、非负，且\(W=\int_a^bw(x)\,dx>0\)，定义\(f\)的{\heiti{}加权平均值}为
\begin{equation}\label{eq:calculus-weighted-average}
 \overline f_w=\frac{\int_a^bf(x)w(x)\,dx}{\int_a^bw(x)\,dx}.
\end{equation}
分子累积各段的加权贡献，分母累积总权重。权为正常数时，就得到普通积分平均值。

从有限分割看，若第\(i\)小段的权为\(w(\xi_i)\Delta x_i\)，离散加权平均就是
\begin{equation*}
 \frac{\sum_i f(\xi_i)w(\xi_i)\Delta x_i}{\sum_iw(\xi_i)\Delta x_i}.
\end{equation*}
分割加细后，分子与分母分别趋于对应积分，便得到式\eqref{eq:calculus-weighted-average}。

若\(m\le f\le M\)，将不等式乘以\(w\ge0\)并积分，得到\(m\le\overline f_w\le M\)。若\(f\)连续，再由介值定理，存在\(\xi\in[a,b]\)，使\(\overline f_w=f(\xi)\)。

例如，线密度为\(\rho\)的细杆可看成许多小质量的组合。其平均位置，即质心坐标，为
\begin{equation*}
 \overline x=\frac{\int_a^bx\rho(x)\,dx}{\int_a^b\rho(x)\,dx}.
\end{equation*}
它由位置乘对应质量的总和除以总质量得到，与有限个质点的质量加权平均一致。

合并两个部分时也应按各自总权重加权。若两段总权重为\(W_1,W_2\)，平均值为\(m_1,m_2\)，则由积分的区间可加性，整体平均值为\((W_1m_1+W_2m_2)/(W_1+W_2)\)。只有两部分总权重相等时，才能对两个平均值直接取算术平均。

\begin{example}
一根细杆占据\([0,L]\)，线密度为\(\rho(x)=k(1+x/L)\)，其中\(k,L>0\)。求总质量与质心位置。再说明同时将所有位置的密度乘以同一个正常数时，质心为什么不变。
\end{example}

\begin{exercise}
\mbox{}\par\nobreak

（1）设\(v(t)=\cos t\)，\(0\le t\le2\pi\)。求位移、路程和平均速度，并与平均速率比较。

（2）杆的线密度满足\(\rho(a+b-x)=\rho(x)\)。用积分换元证明质心位于中点。

（3）设\(w\ge0\)、\(\int_a^bw>0\)。证明\(\int_a^b(f-c)^2w\)在\(c=\overline f_w\)时最小，并解释加权平均与平方偏差之间的关系。

（4）某弹簧的拉力为\(F(x)=kx\)、\(k>0\)。求将其伸长量从\(a\)增加到\(b\)所需的功，其中\(0\le a<b\)。比较从\(0\)到\(a\)与从\(a\)到\(2a\)的功，解释为何相等位移未必对应相等的功。

\end{exercise}

\endgroup
\FloatBarrier
\newpage
\subsection{参数曲线围成的面积*}

直角坐标下的面积可以经换元写成参数形式。若一段曲线表示为\(x=\varphi(t)\)、\(y=\psi(t)\)，且\(x\)沿该段递增，则曲线与横轴之间的有向面积为
\begin{equation*}
 \int y\,dx=\int\psi(t)\varphi'(t)\,dt.
\end{equation*}
若方向反过来，积分相应变号。

以下设闭曲线除首尾相接外不自交，其坐标在有限分段上连续可导，并能拆成横坐标单调的弧段。沿逆时针方向绕区域一周，在各条带上配对上、下边界，累加切片面积公式，得到
\begin{equation*}
 S=-\oint y\,dx.
\end{equation*}
符号\(\oint\)表示沿闭曲线一周，把各参数段的积分相加。逆时针时下边界从左向右、上边界从右向左，因而负号恰好使上、下高度之差为正。

一条带的下、上边界若为\(g,f\)，横坐标范围为\([u,v]\)，沿边界的贡献为
\begin{equation*}
 \int_u^vg(x)\,dx+\int_v^uf(x)\,dx
 =-\int_u^v(f(x)-g(x))\,dx.
\end{equation*}
把各条带相加，得到整个面积；人为切开的边界两侧方向相反，其贡献相消，竖直连接段的\(dx\)为零。这说明闭曲线公式来自切片面积和边界方向的配合。

又因为一周后\(xy\)回到原值，\(\oint d(xy)=\oint(x\,dy+y\,dx)=0\)，所以也有
\begin{equation}\label{eq:calculus-parametric-area}
 S=\oint x\,dy=\frac12\oint(x\,dy-y\,dx)
 =\frac12\int_\alpha^\beta(\varphi\psi'-\psi\varphi')\,dt.
\end{equation}
本式要求沿边界逆时针走一周。反向走得到负的有向面积，多走几周则重复累积。

对椭圆\(x=a\cos t\)、\(y=b\sin t\)，\(0\le t\le2\pi\)，有\(xy'-yx'=ab\)，故面积为\(\pi ab\)。参数表示把四段边界统一在同一个积分中。

\begin{exercise}
\mbox{}\par\nobreak

（1）对上述椭圆分别取\(t\in[0,4\pi]\)及反向参数，判断式\eqref{eq:calculus-parametric-area}的结果，解释覆盖次数与方向的作用。

（2）设\(r(t)\)表示一条正向简单闭曲线。证明平移后面积不变；再用\(x=au+cv\)、\(y=bu+dv\)代入面积公式，核对有向面积乘以\(ad-bc\)的规律。
\end{exercise}

\subsection{极坐标下的面积*}

设\(r=r(\theta)\ge0\)连续，\(\alpha\le\theta\le\beta\)，其中\(\beta-\alpha\le2\pi\)。研究每条射线上从原点到\(r(\theta)\)之间的区域。把角度分成许多小段，在每段用最小、最大半径作内外扇形。

弧度制下，半径为\(r\)、圆心角为\(\Delta\theta\)的扇形面积是\(r^2\Delta\theta/2\)。因此内外扇形和分别是\(r^2/2\)的下和与上和。连续性使两者趋于同一积分，得到
\begin{equation}\label{eq:calculus-polar-area}
 S=\frac12\int_\alpha^\beta r(\theta)^2\,d\theta.
\end{equation}
角度区间决定覆盖的射线，半径决定每条射线上的长度。若参数重复描绘同一区域，应先选出恰好覆盖一次的范围。

\begin{figure}[!ht]
\centering
\includegraphics[width=.47\linewidth]{calculus-polar-area}
\caption{用取样半径作小扇形，局部面积近似为\(r^2\Delta\theta/2\)，再沿角度累积。}
\label{fig:calculus-polar-area}
\end{figure}

当\(r\)连续可导时，可用参数曲线公式核对。令\(x=r\cos\theta\)、\(y=r\sin\theta\)，求导得
\begin{equation*}
 xy'-yx'=r^2(\cos^2\theta+\sin^2\theta)=r^2.
\end{equation*}
两条连接原点的径向边界不贡献\(x\,dy-y\,dx\)，所以参数面积公式也给出式\eqref{eq:calculus-polar-area}。扇形和仅要求半径函数连续，参数求导则联系了更光滑情形下的两种面积表示。

\begin{example}
求心形线\(r=a(1+\cos\theta)\)、\(0\le\theta\le2\pi\)围成的面积。再用\(r=2a\cos\theta\)、\(-\pi/2\le\theta\le\pi/2\)计算圆的面积，并与直角坐标下的圆方程核对。
\end{example}

\begin{exercise}
\mbox{}\par\nobreak

（1）求曲线\(r=a\cos2\theta\)在\(-\pi/4\le\theta\le\pi/4\)上所围成一瓣的面积，再由对称性求全部四瓣的总面积。

（2）若区域在每条射线上介于\(r_1(\theta)\)与\(r_2(\theta)\)之间，且\(0\le r_1\le r_2\)，推导其面积公式。用同心圆环核对。
\end{exercise}

\subsection{体积*}

对本节研究的几何体，可以把沿某个方向的薄层近似为柱体。设在位置\(x\)处垂直于横轴的截面积为连续函数\(A(x)\)，薄层厚度为\(\Delta x\)，则局部体积近似为\(A(x)\Delta x\)。由薄层和的极限，得到
\begin{equation}\label{eq:calculus-volume}
 V=\int_a^b A(x)\,dx.
\end{equation}
需要先确定截面，再计算它的面积；截面形状不同，对应的\(A(x)\)也不同。

\paragraph{旋转体的截面积} \mbox{}

若将\(0\le y\le f(x)\)、\(a\le x\le b\)的区域绕横轴旋转，垂直于旋转轴的截面为半径\(f(x)\)的圆盘，因此
\begin{equation*}
 V=\pi\int_a^bf(x)^2\,dx.
\end{equation*}
若\(0\le g(x)\le f(x)\)，旋转的是两曲线之间的区域，则截面为圆环，外半径、内半径分别为\(f(x),g(x)\)，故
\begin{equation}\label{eq:calculus-washer-volume}
 V=\pi\int_a^b\bigl(f(x)^2-g(x)^2\bigr)\,dx.
\end{equation}
圆环面积是两个圆面积之差，不能把\(f-g\)当作半径后再平方。若区域跨过旋转轴，则应从旋转后的实际截面重新判断半径。

绕纵轴旋转时，改用水平截面更直接：若在\(c\le y\le d\)上，区域介于\(0\le u(y)\le x\le v(y)\)之间，则
\begin{equation*}
 V=\pi\int_c^d\bigl(v(y)^2-u(y)^2\bigr)\,dy.
\end{equation*}
半径始终是到旋转轴的距离。绕\(y=k\)或\(x=k\)旋转，只须根据位置关系重新写出内、外半径；累积变量仍沿轴的方向选取。

例如，半径为\(R\)的球在横坐标\(x\)处截面积为\(\pi(R^2-x^2)\)，所以
\begin{equation*}
 V=\pi\int_{-R}^R(R^2-x^2)\,dx=\frac43\pi R^3.
\end{equation*}
熟悉的体积公式由截面的累积得到。

截面不必是圆。例如，椭球\(x^2/a^2+y^2/b^2+z^2/c^2\le1\)在横坐标\(x\)处的截面为椭圆，两条半轴分别为\(b\sqrt{1-x^2/a^2}\)、\(c\sqrt{1-x^2/a^2}\)。它的截面积是\(\pi bc(1-x^2/a^2)\)，再沿\([-a,a]\)累积，就将椭球体积化成了一个多项式积分。圆与椭圆的截面虽不同，确定截面、求面积、沿厚度累积的次序相同。

\begin{example}
求\(y=x\)与\(y=x^2\)在\([0,1]\)上围成区域绕横轴旋转所得几何体的体积。再用截面积法求底面半径为\(R\)、高为\(h\)的正圆锥体积。
\end{example}

\paragraph{圆柱壳法} \mbox{}

同一个旋转体也可以按到轴的距离分层。先设\(0\le a<b\)、\(f\ge0\)，将\(a\le x\le b\)、\(0\le y\le f(x)\)的区域绕纵轴旋转。竖直窄条旋转后形成薄圆柱壳，其半径接近\(x\)，高度接近\(f(x)\)，厚度为\(\Delta x\)。

对固定高度\(H\)，内、外半径为\(x_{i-1},x_i\)的圆柱壳体积恰为
\begin{equation*}
 \pi(x_i^2-x_{i-1}^2)H
 =2\pi\frac{x_i+x_{i-1}}2H\Delta x_i.
\end{equation*}
在每段分别以\(f\)的最小值、最大值取代\(H\)，便得到内外近似体积。连续性使高度的波动统一趋零，而圆环底面积之和为固定的\(\pi(b^2-a^2)\)，所以两种近似之差趋零，得到
\begin{equation}\label{eq:calculus-shell-volume}
 V=2\pi\int_a^bxf(x)\,dx.
\end{equation}
若旋转的是\(g(x)\le y\le f(x)\)之间的区域，壳高改为\(f-g\)，于是\(V=2\pi\int_a^bx(f-g)\,dx\)。这里区域位于轴的一侧，各半径层只累积一次。

截面积法沿轴累积圆盘或圆环，圆柱壳法垂直于轴累积半径层。两者可以处理同一个几何体，选择哪一种，取决于边界表达和需要拆分的次数。

\begin{figure}[!ht]
\centering
\includegraphics[width=.83\linewidth]{calculus-rotation}
\caption{左图用圆盘截面沿轴累积，右图的竖直窄条绕纵轴生成圆柱壳；半径、截面积与壳高分别由几何关系确定。}
\label{fig:calculus-rotation}
\end{figure}

\begin{example}
将\(y=x\)与\(y=x^2\)围成的区域绕纵轴旋转，分别用圆环截面与圆柱壳计算体积。再将它绕\(x=-1\)旋转，说明圆柱壳的半径与高度分别怎样变化。
\end{example}

\begin{exercise}
\mbox{}\par\nobreak

（1）将\(y=\sqrt{x}\)、横轴及直线\(x=1\)、\(x=4\)围成的区域绕横轴旋转，求体积。再改绕\(y=-1\)旋转，写出圆环的内、外半径。

（2）将\(y=x^2\)、\(y=1\)围成的区域绕纵轴旋转，分别用水平截面和圆柱壳计算。为什么使用圆柱壳时，不能把左右两半区域各自旋转所得体积直接相加？

（3）补完上文椭球体积的计算，核对结果在\(a=b=c\)时是否回到球体积。
\end{exercise}

\subsection{弧长*}

对平面参数曲线\(r(t)=(x(t),y(t))\)，在参数区间\([\alpha,\beta]\)上作分割\(P\)，将曲线上对应的点依次连接，得到折线长度
\begin{equation*}
 L(P)=\sum_i|r(t_i)-r(t_{i-1})|.
\end{equation*}
若当\(\|P\|\to0\)时，这些长度趋于同一个有限极限，就称该极限为这段曲线的{\heiti{}弧长}。

现在设两个坐标函数连续可导。由局部一次近似，短弦长度应接近
\begin{equation*}
 |r'(t)|\Delta t=\sqrt{x'(t)^2+y'(t)^2}\,\Delta t.
\end{equation*}
把它累积，就得到弧长公式
\begin{equation}\label{eq:calculus-arc-length}
 L=\int_\alpha^\beta\sqrt{x'(t)^2+y'(t)^2}\,dt.
\end{equation}

为说明误差如何控制，对每段分别使用坐标函数的中值定理，短弦增量可写为\(\Delta t_i(x'(\xi_i),y'(\eta_i))\)。两个中间点未必相同，但导函数的一致连续性保证，当段长足够小时，这个向量与\(\Delta t_i r'(\tau_i)\)之差的长度不超过\(\varepsilon\Delta t_i\)。由向量长度的三角不等式，弦长与\(|r'(\tau_i)|\Delta t_i\)的差也受此界控制。累积后总误差不超过\(\varepsilon(\beta-\alpha)\)，故折线长度与速度模的Riemann和趋于同一极限，证明了公式。

这里沿用第二章的一致连续性，是将逐点的线性近似用于整个曲线长度的依据。对图像\(y=f(x)\)，取参数为\(x\)，便得到
\begin{equation*}
 L=\int_a^b\sqrt{1+f'(x)^2}\,dx.
\end{equation*}
几何长度累积的是速度的模，而不是某一个坐标变化率。

\begin{figure}[!ht]
\centering
\includegraphics[width=.83\linewidth]{calculus-arc-length}
\caption{左图用折线近似弧长；右图将短弦的两个坐标增量合成为长度，再过渡到弧长积分。}
\label{fig:calculus-arc-length}
\end{figure}

若曲线更适合写成\(x=g(y)\)，则取\(y\)为参数，得到\(L=\int_c^d\sqrt{1+g'(y)^2}\,dy\)。同一段曲线可以有不同的坐标表示，弧长公式始终累积两个坐标变化合成的长度。

\paragraph{弧长函数} \mbox{}

固定参数起点，令
\begin{equation*}
 s(t)=\int_\alpha^t\sqrt{x'(u)^2+y'(u)^2}\,du.
\end{equation*}
它表示沿参数递增方向，从起点到当前位置已经走过的长度。由基本定理，
\begin{equation*}
 s'(t)=\sqrt{x'(t)^2+y'(t)^2},\qquad
 ds=\sqrt{x'(t)^2+y'(t)^2}\,dt.
\end{equation*}
若\(t\)是时间，右端的导数就是速率。这个长度元也使“沿曲线按长度累积”获得具体表达，例如非均匀细线的质量可以写成\(\int\rho\,ds\)，其中\(\rho\)为单位弧长的质量，计算时再代入参数形式。

参数改变时，走过同一位置所需的参数增量会变化，速度模也相应变化。若\(t=\varphi(s)\)、\(\varphi'>0\)，则新的速度模为\(|r'(\varphi(s))|\varphi'(s)\)。换元积分使它沿新参数的累积仍等于原来的弧长。面积公式保留方向，弧长公式使用长度，这解释了反向走同一段曲线时面积积分变号、长度却不变。

\begin{example}
用参数公式求半径为\(R\)的圆周长。再求曲线\(x=t^2\)、\(y=t^3\)、\(0\le t\le1\)的弧长，比较直接使用参数与先消去参数的计算。
\end{example}

\paragraph{极坐标曲线的弧长} \mbox{}

若\(r=r(\theta)\)连续可导，代入\(x=r\cos\theta\)、\(y=r\sin\theta\)，有
\begin{equation*}
 \begin{aligned}
 x'&=r'\cos\theta-r\sin\theta,\\
 y'&=r'\sin\theta+r\cos\theta,
 \end{aligned}
 \qquad x'^2+y'^2=r'^2+r^2.
\end{equation*}
交叉项相消，因而
\begin{equation}\label{eq:calculus-polar-arc-length}
 L=\int_\alpha^\beta\sqrt{r(\theta)^2+r'(\theta)^2}\,d\theta.
\end{equation}
其中\(r'\)记录半径的变化，\(r\)对应角度变化造成的横向位移。即使\(r\)为常数，角度变化仍产生长度，因此弧长不能只累积\(|r'|\)。

取\(r\equiv R\)，公式给出\(L=R(\beta-\alpha)\)，与圆弧的弧度定义一致；对一般曲线，还需要同时计入半径的变化。

\begin{example}
求心形线\(r=a(1+\cos\theta)\)、\(0\le\theta\le2\pi\)的周长，并与前面的面积计算比较。化简根式时注意\(\sqrt{\cos^2(\theta/2)}=|\cos(\theta/2)|\)。
\end{example}

\begin{exercise}
\mbox{}\par\nobreak

（1）求\(x=y^2\)、\(1\le y\le2\)的弧长，比较以\(x\)为变量和以\(y\)为变量的计算。

（2）摆线的一拱由\(x=a(t-\sin t)\)、\(y=a(1-\cos t)\)、\(0\le t\le2\pi\)表示。求其弧长。

（3）对同一段光滑曲线作保持方向的连续可导参数替换，且替换导数为正。用换元公式证明弧长不变，并与参数曲线切线方向的研究联系起来。

（4）一根细线沿四分之一圆周\(x=R\cos t\)、\(y=R\sin t\)、\(0\le t\le\pi/2\)放置，单位弧长的质量为\(\rho(t)=k(1+\sin t)\)。求总质量，并说明为什么不能直接计算\(\int_0^{\pi/2}\rho(t)\,dt\)。

（5）求极坐标曲线\(r=a\theta\)、\(0\le\theta\le1\)的弧长。把结果与径向距离的总增量比较。
\end{exercise}

\subsection{曲率*}
\label{subsec:calculus-curvature}

第四章用切线和法线研究二次曲线，一阶导数进一步给出了切线方向。不同曲线在一点可以有同一条切线，却以不同的程度离开它。若要比较这种弯曲，就需要研究切线方向沿曲线怎样改变。

同一段曲线可以快走，也可以慢走，参数变化的快慢并不改变曲线本身。因此先用已经建立的弧长衡量走过的距离，再比较单位弧长内切线方向的变化。

\begin{definition}
设平面曲线\(r(t)=(x(t),y(t))\)的坐标二阶连续可导，且\(r'(t)\ne0\)。记速率\(v(t)=|r'(t)|\)，弧长函数为\(s(t)\)，则\(ds/dt=v(t)>0\)。定义单位切向量
\begin{equation*}
 T(t)=\frac{r'(t)}{v(t)},
\end{equation*}
称\(\kappa=|dT/ds|\)为该点的{\heiti{}曲率}。若\(\kappa>0\)，称\(\rho=1/\kappa\)为{\heiti{}曲率半径}。
\end{definition}

局部选取切线方向角\(\theta\)，使\(T=(\cos\theta,\sin\theta)\)，则\(|dT/ds|=|d\theta/ds|\)。所以曲率也就是切线方向角对弧长的变化率的绝对值。直线的切向不变，曲率为零；半径\(R\)的圆走过弧长\(R\Delta\theta\)时，切线方向也转过\(\Delta\theta\)，故曲率恒为\(1/R\)。

\paragraph{参数公式} \mbox{}

由\(x'=v\cos\theta\)、\(y'=v\sin\theta\)，对参数求导并消去\(v'\)，得到
\begin{equation*}
 x'y''-y'x''=v^2\theta'.
\end{equation*}
又因为\(ds/dt=v\)，于是
\begin{equation}\label{eq:calculus-parametric-curvature}
 \kappa=\frac{|\theta'|}{v}
 =\frac{|x'y''-y'x''|}{(x'^2+y'^2)^{3/2}}.
\end{equation}
分子是第三章中两个平面向量\(r',r''\)的行列式的绝对值；它记录速度变化中使方向转动的部分。分母消除了参数速度的影响。反向走同一曲线时，方向角变化的符号可以改变，曲率的大小仍相同。

平移不改变坐标导数，旋转则保持向量长度和行列式，因而这两种变换都保持曲率。第四章中改变坐标原点或旋转坐标轴，不会改变所研究曲线的弯曲程度。

对图像\(y=f(x)\)，取\(x\)为参数，得到
\begin{equation}\label{eq:calculus-graph-curvature}
 \kappa=\frac{|f''(x)|}{(1+f'(x)^2)^{3/2}}.
\end{equation}
因此，\(f''\)的大小本身并不能直接充当弯曲程度；还要计入斜率对弧长的影响。在水平切线处，公式才简化为\(\kappa=|f''|\)。

\paragraph{密切圆} \mbox{}

在曲率非零的点\(P\)，令\(N=(dT/ds)/\kappa\)。由\(T\cdot T=1\)求导，知\(T\cdot N=0\)，且\(|N|=1\)。它指向切线方向正在转向的一侧。取
\begin{equation*}
 C=P+\rho N,
\end{equation*}
以\(C\)为圆心、\(\rho\)为半径作圆，称为曲线在\(P\)处的{\heiti{}密切圆}，也称{\heiti{}曲率圆}。

圆心位于法线上，半径也有确定的选择。以\(P\)处为弧长起点，曲线的二阶Taylor展开为
\begin{equation*}
 r(s_0+h)=P+Th+\frac12\kappa Nh^2+o(h^2).
\end{equation*}
上述圆可在\(P\)附近写为
\begin{equation*}
 c(h)=P+\rho\sin(h/\rho)T+\rho\bigl(1-\cos(h/\rho)\bigr)N.
\end{equation*}
展开后同样得到\(P+Th+\kappa Nh^2/2+o(h^2)\)。切线匹配了一阶方向，密切圆还匹配二阶弯曲；这使第四章的圆与法线，成为一般曲线局部近似的工具。

\paragraph{圆锥曲线中的曲率} \mbox{}

对椭圆\(x=a\cos t\)、\(y=b\sin t\)，其中\(a>b>0\)，有\(x'y''-y'x''=ab\)，代入式\eqref{eq:calculus-parametric-curvature}得
\begin{equation*}
 \kappa(t)=\frac{ab}{(a^2\sin^2t+b^2\cos^2t)^{3/2}}.
\end{equation*}
长轴端点的曲率半径为\(b^2/a\)，短轴端点则为\(a^2/b\)。前者弯曲较急，后者较平缓。把圆沿两个方向作不同倍率的伸缩，虽然仍保持切线和接触关系，却不保持圆形与曲率；曲率还依赖长度和角度的度量。

对抛物线\(y^2=2px\)、\(p>0\)，沿用第四章的参数\(x=pt^2/2\)、\(y=pt\)，得到
\begin{equation*}
 \kappa(t)=\frac1{p(1+t^2)^{3/2}}.
\end{equation*}
顶点处的曲率半径为\(p\)，密切圆圆心为\((p,0)\)。它与焦点\((p/2,0)\)位于同一条轴上，但由不同的几何关系确定：焦点参与整条曲线的定义，密切圆描述顶点处的局部弯曲。

\begin{figure}[!ht]
\centering
\includegraphics[width=.83\linewidth]{calculus-curvature}
\caption{椭圆长轴端点及抛物线顶点处的密切圆。圆心位于法线上，圆与曲线在接触点匹配切向和曲率。}
\label{fig:calculus-curvature}
\end{figure}

\begin{example}
对椭圆\(x^2/4+y^2=1\)，求长、短轴端点处的密切圆方程。再求抛物线\(y^2=4x\)顶点处的密切圆，比较圆心与焦点的位置。
\end{example}

\begin{exercise}
\mbox{}\par\nobreak

（1）将圆\(x=R\cos t\)、\(y=R\sin t\)代入参数曲率公式，核对\(\kappa=1/R\)。再求\(y=x^2\)在原点及\(x=1\)处的曲率，比较它与常值二阶导数\(2\)。

（2）用\(x=a\cosh t\)、\(y=b\sinh t\)求双曲线\(x^2/a^2-y^2/b^2=1\)的曲率。求右顶点处密切圆的圆心与半径，并与椭圆长轴端点比较。

（3）对保持方向的二阶连续可导参数替换\(t=\varphi(u)\)、\(\varphi'>0\)，由式\eqref{eq:calculus-parametric-curvature}证明曲率不变。二阶求导产生的哪一项在行列式中消失？

（4）将平面图形整体放大\(c>0\)倍，证明曲率变为原来的\(1/c\)，曲率半径变为原来的\(c\)倍。再用圆变椭圆的例子，说明第四章的一般仿射变换为什么没有同一个曲率倍率。

（5）椭圆\(x^2/a^2+y^2/b^2=1\)、\(a>b>0\)在哪些点弯曲最急，哪些点最平缓？求两处曲率的比值，并讨论当\(a/b\)增大时怎样变化。
\end{exercise}

\subsection{旋转曲面的面积*}

曲线绕直线旋转时，每个点扫过一个圆，整条母线生成一片曲面。旋转体体积累积的是截面积乘厚度，曲面面积则需要同时考虑旋转半径与母线的长度。以下选取位于旋转轴一侧、恰好生成曲面一次的母线。

\paragraph{从圆台侧面积到曲面面积} \mbox{}

先设\(f\ge0\)且连续可导，将\(y=f(x)\)、\(a\le x\le b\)绕横轴旋转。用母线的折线近似曲线，每条线段旋转后生成圆台侧面；若两端半径相同，就是圆柱侧面。

第\(i\)段母线长为
\begin{equation*}
 \ell_i=\sqrt{(\Delta x_i)^2+(f(x_i)-f(x_{i-1}))^2},
\end{equation*}
由圆台侧面积公式，该段旋转后的面积为\(\pi(f(x_{i-1})+f(x_i))\ell_i\)。对本节的光滑旋转曲面，由这些圆台侧面积和在分割加细时的极限确定面积。

\begin{figure}[!ht]
\centering
\includegraphics[width=.48\linewidth]{calculus-surface-frustum}
\caption{短弦绕轴生成圆台侧面，其面积为两端圆周长的平均值乘短弦长度；加细后得到\(2\pi r\,ds\)的累积。}
\label{fig:calculus-surface-frustum}
\end{figure}

对\(f\)使用中值定理，可写\(\ell_i=\sqrt{1+f'(\eta_i)^2}\Delta x_i\)。连续性使两端半径的平均与\(f(\eta_i)\)的差一致趋零。又因\(f'\)有界，设\(|f'|\le K\)，则\(\sum_i\ell_i\le\sqrt{1+K^2}(b-a)\)，所以替换半径平均造成的总误差也趋零。圆台和因此与下列Riemann和有相同极限，得到
\begin{equation}\label{eq:calculus-surface-x}
 S=2\pi\int_a^bf(x)\sqrt{1+f'(x)^2}\,dx.
\end{equation}
其中\(2\pi f(x)\)是旋转一周的周长，\(\sqrt{1+f'(x)^2}\,dx\)是母线的长度元。直接写成\(2\pi\int f(x)\,dx\)，只计入横向宽度，会漏掉母线倾斜造成的长度。

令\(f\equiv R\)，公式得到圆柱侧面积\(2\pi R(b-a)\)。若母线从半径\(r_1\)线性变化到\(r_2\)，轴向长度为\(h\)，则斜率恒定，公式给出
\begin{equation*}
 S=\pi(r_1+r_2)\sqrt{h^2+(r_2-r_1)^2},
\end{equation*}
与圆台的几何公式一致。

\paragraph{参数表示与旋转轴} \mbox{}

若母线用\(x=x(t)\)、\(y=y(t)\)表示，坐标连续可导，且\(y\ge0\)，绕横轴旋转时同样以短弦构成圆台。弧长证明中的一致误差控制，加上半径函数的连续性，得到
\begin{equation}\label{eq:calculus-parametric-surface}
 S_x=2\pi\int_\alpha^\beta y(t)\sqrt{x'(t)^2+y'(t)^2}\,dt.
\end{equation}
母线可以先分成有限段沿旋转轴坐标单调的弧段，再把各段面积相加，参数不必在整段上都能反解成一个横坐标函数。

绕纵轴旋转、取\(x\ge0\)的母线时，旋转半径改为\(x(t)\)，所以
\begin{equation*}
 S_y=2\pi\int_\alpha^\beta x(t)\sqrt{x'(t)^2+y'(t)^2}\,dt.
\end{equation*}
更一般地，若母线上各点到所选旋转轴的距离为\(\rho(t)\)，面积统一由\(2\pi\rho\,ds\)累积。改变旋转轴只改变半径，长度元仍由同一条母线决定。

这里计算母线扫出的曲面。若求封闭几何体的总表面积，还须加上未由母线生成的底面或端面；若重复扫过同一片曲面，应先选取只覆盖一次的范围。

\begin{example}
将上半圆\(x=R\cos t\)、\(y=R\sin t\)、\(0\le t\le\pi\)绕横轴旋转，求球面积。再求其中两平行平面\(x=c_1\)、\(x=c_2\)之间的球带面积，观察它与球带轴向高度的关系。
\end{example}

\begin{example}
将\(y=\sqrt{x}\)、\(1\le x\le4\)绕横轴旋转，求生成的曲面面积。再写出该曲线绕直线\(y=-1\)旋转时的面积积分，说明哪些因子改变、哪些因子保持不变。
\end{example}

\begin{exercise}
\mbox{}\par\nobreak

（1）用积分求底面半径为\(R\)、高为\(h\)的正圆锥侧面积，再求总表面积。比较体积与侧面积中局部贡献的表达式。

（2）将\(y=x\)、\(1\le x\le2\)分别绕横轴、纵轴旋转，求曲面面积，用圆台公式核对。

（3）设\(R>r>0\)。圆\(x=R+r\cos t\)、\(y=r\sin t\)、\(0\le t\le2\pi\)绕纵轴旋转生成圆环面，求其面积。解释为何此处应选整圆作为母线，而求球面时只选一个半圆。

（4）同一条母线改变参数后，利用换元公式证明旋转面积不变。若沿母线走两遍，积分结果如何变化？
\end{exercise}
